\documentclass[../main.tex]{subfiles}
\graphicspath{{\subfix{../images/}}}
\begin{document}
\newpage
\addcontentsline{toc}{section}{Introduction Générale}
\section*{Introduction Générale}
En Tunisie, plusieurs restrictions sur le transfert d'argent en ligne existent et limitent la liberté du Citoyen Tunisien dans son interaction avec le web en générale. 

Ces limitations restreignent notre communication avec le monde externe mais également les opportunités d'affaire dans le monde interne. Plusieurs changements sont en marche mais ce n'est pas suffisant. Malgré tout, le pays ne se démocratise pas assez rapidement pour offrir à nos entrepreneurs assez d'options pour rester et investir dans le pays. 

C'est dans ce cadre que AYV Studio et Watto ont décider d'entreprendre et de travailler ensemble pour offrir un peu plus de possibilités et d'options au Tunisien quand il s'agit d'acheter en ligne. En effet, dans l'optique d'offrir à notre pays une option de plus. Watto a décidé d'investir sur un projet de cagnottes dont l'unique but est de permettre à tout citoyen d'investir dans n'importe quel cagnotte quel qu'on soit le but (tant que cela reste dans le cadre de la légalité). 

\vspace{10pt}
\newline
C'est là qu'interviens, AYV Studio une plateforme de lancement et de réalisation de tout projet digital. Le partenariat entre les deux entreprises a fait de ce projet ce qu'il est maintenant 

\newline Ce rapport présentera les étapes de réalisation du projet en 7 Chapitres:
\begin{itemize}
  \item \textbf{Chapitre 1}: Cadre Générale, ce chapitre présentera les deux parties prenantes du projet ainsi que la revue de la littérature et les objectifs du projet.
    \item \textbf{Chapitre 2}: Approche de développement et Technologies utilisées, ce chapitre abordera en profendeur notre méthodologie de travail ainsi les technologies et outils utilisées.
    \item \textbf{Chapitre 3}: Spécification des éxigences, ce chapitre décris de façon symbolique et schématique les diagrammes des différents éléments et composants du système.
    \item \textbf{Chapitre 4}: Itération 1
    \item \textbf{Chapitre 5}: Itération 2
    \item \textbf{Chapitre 6}: Itération 3
    \item \textbf{Chapitre 7}: Conclusion générale
\end{itemize}

\end{document}